\documentclass[10pt,a4paper]{amsart}
\usepackage{amsmath,amssymb,amsthm,graphicx}
\usepackage{verbatim}
\newenvironment{reference}{\comment}{\endcomment}
%\newenvironment{reference}{}{}
\newenvironment{slogan}{\comment}{\endcomment}
\newenvironment{history}{\comment}{\endcomment}

% For commutative diagrams we use Xy-pic
\usepackage[all]{xy}

% We use 2cell for 2-commutative diagrams.
\xyoption{2cell}
% Theorem environments.
%
\theoremstyle{plain}
\newtheorem{theorem}[subsection]{Theorem}
\newtheorem{proposition}[subsection]{Proposition}
\newtheorem{lemma}[subsection]{Lemma}

\theoremstyle{definition}
\newtheorem{definition}[subsection]{Definition}
\newtheorem{example}[subsection]{Example}
\newtheorem{exercise}[subsection]{Exercise}
\newtheorem{situation}[subsection]{Situation}

\theoremstyle{remark}
\newtheorem{remark}[subsection]{Remark}
\newtheorem{remarks}[subsection]{Remarks}

\numberwithin{equation}{subsection}

% Macros
%
\def\lim{\mathop{\mathrm{lim}}\nolimits}
\def\colim{\mathop{\mathrm{colim}}\nolimits}
\def\Spec{\mathop{\mathrm{Spec}}}
\def\Hom{\mathop{\mathrm{Hom}}\nolimits}
\def\Ext{\mathop{\mathrm{Ext}}\nolimits}
\def\SheafHom{\mathop{\mathcal{H}\!\mathit{om}}\nolimits}
\def\SheafExt{\mathop{\mathcal{E}\!\mathit{xt}}\nolimits}
\def\Sch{\mathit{Sch}}
\def\Mor{\mathop{\mathrm{Mor}}\nolimits}
\def\Ob{\mathop{\mathrm{Ob}}\nolimits}
\def\Sh{\mathop{\mathit{Sh}}\nolimits}
\def\NL{\mathop{N\!L}\nolimits}
\def\CH{\mathop{\mathrm{CH}}\nolimits}
\def\proetale{{pro\text{-}\acute{e}tale}}
\def\etale{{\acute{e}tale}}
\def\QCoh{\mathit{QCoh}}
\def\Ker{\mathop{\mathrm{Ker}}}
\def\Im{\mathop{\mathrm{Im}}}
\def\Coker{\mathop{\mathrm{Coker}}}
\def\Coim{\mathop{\mathrm{Coim}}}

% Boxtimes
%
\DeclareMathSymbol{\boxtimes}{\mathbin}{AMSa}{"02}

%
% Macros for moduli stacks/spaces
%
\def\QCohstack{\mathcal{QC}\!\mathit{oh}}
\def\Cohstack{\mathcal{C}\!\mathit{oh}}
\def\Spacesstack{\mathcal{S}\!\mathit{paces}}
\def\Quotfunctor{\mathrm{Quot}}
\def\Hilbfunctor{\mathrm{Hilb}}
\def\Curvesstack{\mathcal{C}\!\mathit{urves}}
\def\Polarizedstack{\mathcal{P}\!\mathit{olarized}}
\def\Complexesstack{\mathcal{C}\!\mathit{omplexes}}
% \Pic is the operator that assigns to X its picard group, usage \Pic(X)
% \Picardstack_{X/B} denotes the Picard stack of X over B
% \Picardfunctor_{X/B} denotes the Picard functor of X over B
\def\Pic{\mathop{\mathrm{Pic}}\nolimits}
\def\Picardstack{\mathcal{P}\!\mathit{ic}}
\def\Picardfunctor{\mathrm{Pic}}
\def\Deformationcategory{\mathcal{D}\!\mathit{ef}}
\usepackage{cleveref}
\crefname{originaltheorem}{original theorem}{original theorems}
\Crefname{originaltheorem}{Original theorem}{Original theorems}
\makeatletter
\newlabel{etale-theorem-etale-radicial-open}{{025G}{source}{}{sourcepointer-1}{}}
\newlabel{etale-theorem-etale-radicial-open@cref}{{[originaltheorem][0][]025G}{[1][0][]source}{}{}{}}
\makeatother
\crefname{originallemma}{original lemma}{original lemmas}
\Crefname{originallemma}{Original lemma}{Original lemmas}
\makeatletter
\newlabel{lemma-base-change-etale}{{0466}{source}{}{sourcepointer-2}{}}
\newlabel{lemma-base-change-etale@cref}{{[originallemma][0][]0466}{[1][0][]source}{}{}{}}
\makeatother
\crefname{originallemma}{original lemma}{original lemmas}
\Crefname{originallemma}{Original lemma}{Original lemmas}
\makeatletter
\newlabel{lemma-base-change-universally-injective}{{03MW}{source}{}{sourcepointer-3}{}}
\newlabel{lemma-base-change-universally-injective@cref}{{[originallemma][0][]03MW}{[1][0][]source}{}{}{}}
\makeatother
\crefname{originallemma}{original lemma}{original lemmas}
\Crefname{originallemma}{Original lemma}{Original lemmas}
\makeatletter
\newlabel{lemma-universally-injective}{{040X}{source}{}{sourcepointer-4}{}}
\newlabel{lemma-universally-injective@cref}{{[originallemma][0][]040X}{[1][0][]source}{}{}{}}
\makeatother
\crefname{originalproposition}{original proposition}{original propositions}
\Crefname{originalproposition}{Original proposition}{Original propositions}
\makeatletter
\newlabel{proposition-locally-quasi-finite-separated-over-scheme}{{03XX}{source}{}{sourcepointer-5}{}}
\newlabel{proposition-locally-quasi-finite-separated-over-scheme@cref}{{[originalproposition][0][]03XX}{[1][0][]source}{}{}{}}
\makeatother
\crefname{originallemma}{original lemma}{original lemmas}
\Crefname{originallemma}{Original lemma}{Original lemmas}
\makeatletter
\newlabel{spaces-lemma-surjective-flat-locally-finite-presentation}{{05VM}{source}{}{sourcepointer-6}{}}
\newlabel{spaces-lemma-surjective-flat-locally-finite-presentation@cref}{{[originallemma][0][]05VM}{[1][0][]source}{}{}{}}
\makeatother
\title{A universally injective etale morphism of algebraic spaces is an open immersion}
\author{The Stacks Project Authors}
\date{stacks-project:spaces-morphisms.tex:local-snapshot-20261012; GNU FDL 1.2 or later}
\begin{document}
\maketitle
\setcounter{subsection}{1}
\noindent\textit{Editorial note.} All delivered mathematical text is the author's own, copied from the pinned Stacks Project chapter file spaces-morphisms.tex; no statement or proof step was rewritten, reordered or removed.\par\medskip
\noindent\textit{Editorial note.} The statement begins with a slogan block (a chapter-internal display environment): it is part of the counted statement span, verbatim, and its definition is covered by the reference/slogan/history preamble span, so the delivered text typesets the statement without the slogan line.\par\medskip
\noindent\textit{Editorial note.} The six references of the counted spans are recorded as non-load-bearing navigation pointers; the argument that consumes them stays inside the counted proof span.\par\medskip
\noindent\textit{Editorial note.} Preamble: the three native excerpt spans of sources/stacks-spaces-morphisms/preamble.tex (lines 8-13, 14-18 and 154-217). The chapter file itself carries no preamble text.\par\medskip
\noindent\textit{Allowed standard prerequisites (editorial).}\begin{itemize}
\item Stability of etaleness and universal injectivity under base change. Lemmas lemma-base-change-etale and lemma-base-change-universally-injective, used at proof lines 10958-10962. Reduces the claim to the case where Y is a scheme.
\item The diagonal is etale, a monomorphism and surjective for a universally injective morphism. Lemma lemma-universally-injective, used at proof lines 10963-10965. Supplies the three properties of the diagonal.
\item A surjective flat locally finitely presented monomorphism is an isomorphism. Spaces, Lemma spaces-lemma-surjective-flat-locally-finite-presentation, used at proof lines 10966-10968. Concludes that the diagonal is an isomorphism, i.e. X is separated over Y.
\item A separated locally quasi-finite morphism with representable target has representable source. Proposition proposition-locally-quasi-finite-separated-over-scheme, used at proof lines 10969-10971. Makes X a scheme, so that the scheme theorem applies.
\item The fundamental theorem for etale morphisms of schemes: etale radicial equals open immersion. Etale Morphisms, Theorem etale-theorem-etale-radicial-open, used at proof lines 10972-10974. Concludes that f is an open immersion.
\end{itemize}\medskip
\ifdefined\hypertarget\hypertarget{sourcepointer-1}{}\fi
\noindent\textit{Source reference.} Original theorem 025G: focused-seeds/AGP/sources/stacks-snapshot-20261011/tags.txt, source lines 2813--2813. Navigation pointer into the Stacks Project; the standard fact it supplies is applied inside the counted proof span and the argument resting on it stays there.\par\medskip
\ifdefined\hypertarget\hypertarget{sourcepointer-2}{}\fi
\noindent\textit{Source reference.} Original lemma 0466: focused-seeds/AGP/sources/stacks-snapshot-20261011/tags.txt, source lines 5461--5461. Navigation pointer into the Stacks Project; the standard fact it supplies is applied inside the counted proof span and the argument resting on it stays there.\par\medskip
\ifdefined\hypertarget\hypertarget{sourcepointer-3}{}\fi
\noindent\textit{Source reference.} Original lemma 03MW: focused-seeds/AGP/sources/stacks-snapshot-20261011/tags.txt, source lines 4764--4764. Navigation pointer into the Stacks Project; the standard fact it supplies is applied inside the counted proof span and the argument resting on it stays there.\par\medskip
\ifdefined\hypertarget\hypertarget{sourcepointer-4}{}\fi
\noindent\textit{Source reference.} Original lemma 040X: focused-seeds/AGP/sources/stacks-snapshot-20261011/tags.txt, source lines 5272--5272. Navigation pointer into the Stacks Project; the standard fact it supplies is applied inside the counted proof span and the argument resting on it stays there.\par\medskip
\ifdefined\hypertarget\hypertarget{sourcepointer-5}{}\fi
\noindent\textit{Source reference.} Original proposition 03XX: focused-seeds/AGP/sources/stacks-snapshot-20261011/tags.txt, source lines 5164--5164. Navigation pointer into the Stacks Project; the standard fact it supplies is applied inside the counted proof span and the argument resting on it stays there.\par\medskip
\ifdefined\hypertarget\hypertarget{sourcepointer-6}{}\fi
\noindent\textit{Source reference.} Original lemma 05VM: focused-seeds/AGP/sources/stacks-snapshot-20261011/tags.txt, source lines 7556--7556. Navigation pointer into the Stacks Project; the standard fact it supplies is applied inside the counted proof span and the argument resting on it stays there.\par\medskip
\begin{lemma}
\label{lemma-etale-universally-injective-open}
\begin{slogan}
Universally injective \'etale maps are open immersions.
\end{slogan}
Let $S$ be a scheme. Let $f : X \to Y$ be an \'etale and universally
injective morphism of algebraic spaces over $S$. Then $f$ is an open
immersion.
\end{lemma}

\begin{proof}
Let $T \to Y$ be a morphism from a scheme into $Y$.
If we can show that $X \times_Y T \to T$ is an open immersion, then we
are done. Since being \'etale and being universally injective are
properties of morphisms stable under base change (see
Lemmas \ref{lemma-base-change-etale} and
\ref{lemma-base-change-universally-injective})
we may assume that $Y$ is a scheme. Note that the
diagonal $\Delta_{X/Y} : X \to X \times_Y X$ is \'etale, a monomorphism, and
surjective by
Lemma \ref{lemma-universally-injective}.
Hence we see that $\Delta_{X/Y}$ is an isomorphism (see
Spaces, Lemma
\ref{spaces-lemma-surjective-flat-locally-finite-presentation}),
in particular we see that $X$ is separated over $Y$.
It follows that $X$ is a scheme too, by
Proposition \ref{proposition-locally-quasi-finite-separated-over-scheme}.
Finally, $X \to Y$ is an open immersion by the fundamental theorem
for \'etale morphisms of schemes, see
\'Etale Morphisms, Theorem \ref{etale-theorem-etale-radicial-open}.
\end{proof}

\end{document}
